% ===================================================================
%  तक्ता क्रमांक ७ — हिवाळ्यातले गाव. वसंतातले शेतघर.
%
%  Traduction, faite en 2026, en marathi d'aujourd'hui, sur l'ido de
%  texto/io. Regles de la colonne : voir 10-takta-01.tex. Deux
%  scenes, deux numerotations : les cles portent c1 et c2. Le
%  feuillet 49 porte aussi la ligne d'apparat « दुसरी मालिका ».
%
%  LE VIEUX CORDONNIER FAIT EN MARATHI EXACTEMENT LA MEME BLAGUE
%  QU'EN EGYPTIEN, ET C'EST LA TROISIEME COLONNE OU ELLE TOMBE
%  JUSTE. « En ville je m'appelais shuifisto et je n'avais pas de
%  clients ; ici je m'appelle shureparisto et le travail ne manque
%  pas. » L'egyptien opposait le savant صانع أحذية au turc جزمجي ;
%  le marathi oppose le sanskrit पादत्राणनिर्माता au nom de la caste,
%  चांभार. Deux registres, deux langues, une seule chute.
%
%  LE VILLAGE PORTE SES METIERS PAR LEUR NOM DE CASTE, comme au
%  tableau 5 : चांभार le cordonnier, न्हावी le barbier, लोहार le
%  forgeron, धनगर le berger, गडी le valet de ferme. Et le garde
%  champetre est कोतवाल — l'officier du village marathe, en fonction
%  bien avant le gendarme du livret.
%
%  LA FERME EST LA OU LE MARATHI EST LE PLUS RICHE, parce que c'est
%  la que la langue a le plus servi :
%
%    शिवार (3)    la campagne cultivee   तास (8)    le sillon
%    नांगर (5)    la charrue             कुळव (40)  la herse
%    टिकाव (39)   la pioche              वैरण (10)  le foin
%    विहीर (29)   le puits               डोण (26)   l'abreuvoir
%    गव्हाण       le ratelier            गंजी (33)  la meule
%    खळे          l'aire                 उकिरडा (81) le fumier
%    गोठा (58)    l'etable               खुराडे (44) le poulailler
%    मेंढवाडा (48) la bergerie           कळप (49)   le troupeau
%    चरखा         le rouet               रवी (61)   la baratte
%    टाप          le sabot               नाल (6)    le fer
%    ओढा (23)     le ruisseau            शिंगरू (22) le poulain
%
%  ET LE TROUPEAU SE NOMME BETE PAR BETE, sans composer comme l'ido :
%  मेंढा (50) le belier, मेंढी (53) la brebis, कोकरू (54) l'agneau,
%  शेळी (51) la chevre, करडू (52) le chevreau, वळू (27) le taureau,
%  वासरू (24) le veau, घोडी (21) la jument, शिंगरू (22) le poulain.
%
%  UN EMPRUNT ASSUME : le marron chaud du tableau 18 n'a pas de nom
%  marathi, parce que le chataignier ne pousse pas au Maharashtra.
%  On ecrit चेस्टनट et l'on ne transpose pas la chose : la regle de
%  la colonne quebecoise vaut ici aussi — on modernise la langue,
%  jamais les choses. Ecrire भाजके शेंगदाणे serait traduire le
%  Maharashtra, pas le livret.
%
%  TROIS INVERSIONS PREVENUES, CINQ QUI ONT QUAND MEME PASSE. Les
%  trois vues d'avance : la roue arretee par la glace du ruisseau,
%  la litiere de la jument, les outils oublies pres de la niche du
%  chien. Les cinq autres — le tas de neige pres du pont, la charrue
%  attelee au cheval, le foin retourne a la fourche, la clochette de
%  la vache, la poule sur le fumier — sont sorties a l'envers sans
%  que rien se voie, et c'est renvoji.py qui les a trouvees. Savoir
%  ou la postposition va mordre ne dispense pas de verifier : le
%  controle en trouve toujours plus que l'attention.
% ===================================================================

%%K t07-noto-1 noto
\VUnotes{143.2mm}{%
(*) जेवणाच्या तपशिलांसाठी तक्ते ६ आणि १३ पाहा.%
}

%%K t07-apar-1 apar
\VUsaut{4.0mm}
\VUcentre{13.2pt}{90}{दुसरी मालिका}
\VUsaut{3.0mm}
\VUfilet{16mm}
\VUsaut{5.6mm}
\VUtitre{15.0pt}{60}{तक्ता क्रमांक ७}
\VUsaut{3.0mm}

%%K t07-tit-1 sub
\VUcentre{12.0pt}{0}{\textit{पहिला प्रवेश.}}
\VUsaut{3.0mm}

%%K t07-tit-2 sub
\VUpk{10.2pt}{40}{हिवाळ्यातले गाव.}
\VUsaut{4.0mm}

%%K t07-c1-01-1 p
१. --- नववर्षाच्या सुट्टीत मी माझ्या वडिलांबरोबर नवनगरला गेलो — ते एक लहानसे गाव, जिथे त्यांना
काही व्यापारी कामे उरकायची होती.

%%K t07-c1-02-1 p
२. --- शहर आणि ते गाव यांच्यामध्ये चालणारी \VUgras{टपालगाडी} \textsuperscript{(11)} आम्ही
घेतली; पण फार थंडी असल्यामुळे त्या फारशा सोयीच्या नसलेल्या गाडीतला प्रवास काही फार सुखाचा झाला
नाही.

%%K t07-c1-03-1 p
३. --- म्हणूनच \VUgras{गाडीवानाने} \textit{पांढऱ्या अस्वलाच्या} \VUgras{खानावळीसमोरच्या}
\textsuperscript{(2)} चौकात गाडी थांबवलेली पाहून आम्हाला खरा आनंद झाला. \VUgras{खानावळवाला}
\textsuperscript{(4)} आपल्या घराच्या उंबरठ्यावर शांतपणे \VUgras{चिलीम} ओढत आमची वाट पाहत होता.

%%K t07-c1-03-2 p
त्याने आम्हाला आत बोलावले, आणि चुलीत तडतडणाऱ्या वाळक्या फांद्यांच्या विस्तवासमोर बसायला मी
काही आढेवेढे घेतले नाहीत. मग जुनी \VUgras{पाटी} \textsuperscript{(3)} करकरायला लावणाऱ्या आणि
धुराड्यातून निघणारा \VUgras{धूर} \textsuperscript{(1)} काळ्या भोवऱ्यांनी उडवून नेणाऱ्या त्या
बोचऱ्या \VUgras{उत्तरेच्या वाऱ्याची} मी चांगलीच थट्टा केली.

%%K t07-c1-04-1 p
४. --- न्याहारीची वेळ झाली होती. फार वाट न पाहता, आम्हाला वाढलेले साधे पण चवदार जेवण आम्ही
मनसोक्त जेवलो (*).

%%K t07-tit-3 sub
\VUpk{10.2pt}{40}{काही भेटी.}
\VUsaut{3.0mm}

%%K t07-c1-05-1 p
५. --- न्याहारीनंतर मी माझ्या वडिलांबरोबर गेलो — ते विमा उतरवणारे आहेत — त्यांच्या
गिऱ्हाइकांकडे. आधी आम्ही \VUgras{लोहाराकडे} \textsuperscript{(5)} गेलो, जो खानावळीजवळ राहतो.
तो घोड्याला नाल ठोकण्यात गुंतला होता. त्याच्या सोबत्याच्या ताकदीचे मला कौतुक वाटले — त्याने
घोड्याची \VUgras{टाप} आपल्या कातडी अपऱ्यावर घट्ट धरून ठेवली होती, आणि तेवढ्यात लोहाराने जोरदार
घाव घालून \VUgras{नाल} \textsuperscript{(6)} बसवली.

%%K t07-c1-06-1 p
६. --- मग आम्ही गावच्या \VUgras{चांभाराकडे} \textsuperscript{(7)} गेलो, म्हातारा जेकब. पाटी
म्हणून त्याच्याकडे एक सुंदर \VUgras{बूट} असला, तरी भोके पडलेल्या जुन्या जोड्यांना नवे तळ
लावण्यावरच तो समाधानी होता.

%%K t07-c1-06-2 p
तो म्हातारा हसत आम्हाला म्हणाला : « शहरात मी « पादत्राणनिर्माता » होतो आणि माझ्याकडे गिऱ्हाईक
नव्हते. इथे मी « चांभार » आहे आणि कामाला तोटा नाही. »

%%K t07-c1-07-1 p
७. --- त्याने आम्हाला आपल्या शेजारच्या \VUgras{न्हाव्याबद्दल} \textsuperscript{(8)} सांगितले,
ज्याची गिऱ्हाइके नाराज असतात. त्याचे \VUgras{वस्तरे} नेहमी दातेरी झालेले असतात आणि त्याच्या
\VUgras{कात्र्या} कधीच नीट कापत नाहीत. पण तो गावातला एकमेव न्हावी असल्यामुळे, दाढी आणि केस
त्याच्याकडूनच करून घ्यावे लागतात, हे उघड आहे. शिवाय तो कंजूष आणि तुसडा माणूस आहे.

%%K t07-c1-08-1 p
८. --- सुदैवाने बाकीचे \VUgras{शेजारी} त्याच्यासारखे नाहीत. उदाहरणार्थ \VUgras{गाडीकाम करणारा}
\textsuperscript{(9)} भला माणूस आहे, कष्टाळू आणि मदतीला तत्पर. त्याच्याकडून गाडी दुरुस्त करून
घेणे हा आनंद आहे. त्याचे काम नेहमी पुरे आणि नीटनेटके असते.

%%K t07-c1-09-1 p
९. --- म्हाताऱ्या जेकबने \VUgras{खोगीर बनवणाऱ्याबद्दलही} \textsuperscript{(10)} तक्रार केली
नाही; काम निकडीचे असले की तो कधी कधी त्याला \VUgras{घोड्याचे सामान} शिवायला मदत करतो.

%%K t07-c1-09-2 p
माझ्या वडिलांनी त्याला त्याच्या घरच्यांबद्दल विचारले : « सगळे बरे आहेत. माझी नात
\VUgras{शाळेत} \textsuperscript{(12)} आहे. तिची \VUgras{शिक्षिका} \textsuperscript{(13)}
तिच्यावर फार खूश आहे, ती चांगली \VUgras{विद्यार्थिनी} \textsuperscript{(14)} आहे. माझ्या वाया
गेलेल्या मुलासारखी ती नाही; त्याला फक्त खेळण्याचेच वेध लागलेले असतात. \VUgras{शिक्षकाला}
\textsuperscript{(15)} त्याला काहीही शिकवता येत नाही. »

%%K t07-tit-4 sub
\VUsaut{3.0mm}
\VUpk{10.2pt}{40}{गावातल्या गप्पा.}
\VUsaut{3.0mm}

%%K t07-c1-10-1 p
१०. --- मग त्या म्हाताऱ्याने आम्हाला गावच्या बातम्या सांगितल्या. नवी शाळा बांधणार आहेत, कारण
\VUgras{ग्रामपंचायतीच्या घरातली} शाळा फार लहान पडू लागली. शिवाय ते सोपे आहे, कारण
\VUgras{गिरणीवाल्याच्या} बागेचा एक तुकडा विकत घेतला की भागेल. त्या बिचाऱ्या माणसाला त्याचा फार
फायदा होईल. थंडीमुळे \VUgras{पाणचक्कीची} \textsuperscript{(16)} \VUgras{जाती} गहू दळू शकत
नाहीत, कारण \VUgras{चाक} \textsuperscript{(17)} \VUgras{बर्फाने} \textsuperscript{(25)} गोठून
थांबले आहे — \VUgras{ओढ्याच्या} \textsuperscript{(23)}.

%%K t07-c1-11-1 p
११. --- « बांधकामाबद्दल म्हणाल तर, आमचे नवे \VUgras{चर्च} \textsuperscript{(18)} तुम्ही पाहिले
का? बर्फाच्या शालीखाली ते फार चित्रमय दिसते. जुने घंटाघर आता ओळखू न आल्यामुळे \VUgras{कावळे}
\textsuperscript{(19)} \VUgras{दफनभूमीतल्या} \textsuperscript{(20)} मोठ्या झाडावर उदासपणे
बसतात. »

%%K t07-c1-12-1 p
१२. --- दफनभूमीच्या त्या उल्लेखाने चांभाराला माझ्या वडिलांच्या ओळखीची आणि गेल्या वर्षी वारलेली
माणसे आठवली. « केवढी \VUgras{थडगी} \textsuperscript{(21)}, महाशय, \VUgras{सायप्रसखाली}
\textsuperscript{(22)} बाकीच्यांमध्ये भर पडली! मृत्यूने आमचा म्हातारा नोटरी कापून नेला. सगळे
गाव त्याच्या \VUgras{अंत्ययात्रेला} होते. »

%%K t07-c1-13-1 p
१३. --- « हा बघा त्याचा वारसदार, ओढ्यावर घसरत जाणारा. आपल्या \VUgras{स्केटचा}
\textsuperscript{(26)} तुटलेला पट्टा माझ्याकडून दुरुस्त करून घ्यायला तो नुकताच येऊन गेला. »

%%K t07-c1-14-1 p
१४. --- « आणि हा बघा ओढ्याच्या काठावर, नवा \VUgras{कोतवाल} \textsuperscript{(27)}. तो निवृत्त
पोलिस आहे आणि गावच्या खोडकर पोरांना धाक घालतो. त्याच्या \VUgras{पायपट्ट्या}
\textsuperscript{(29)} आणि त्याची \VUgras{टोपी} \textsuperscript{(28)} दिसताच ती लगेच पळ
काढतात. »

%%K t07-c1-15-1 p
१५. --- « अरे! हा बघा म्हातारा योसेफ, पाववाल्यासाठी \VUgras{सरपणाच्या मोळ्यांची गाडी}
\textsuperscript{(30)} हाकत नेणारा. --- खरेच, माझे वडील म्हणाले, मी त्याची \VUgras{खेचरांची}
\textsuperscript{(31)} जोडी ओळखतो. मला त्याच्याशी बोलायचेच आहे, ही संधी मी साधतो. येतो,
जेकबबाबा. »

%%K t07-c1-16-1 p
१६. --- आम्ही बाहेर पडलो तेवढ्यात गावच्या मुलांनी शाळा सोडली आणि धावत \VUgras{घसरगुंडीकडे}
\textsuperscript{(33)} धाव घेतली — \VUgras{पडून} \textsuperscript{(34)} हातपाय मोडून घेण्याचा
धोका पत्करून. त्यांच्यातल्या एकाने गाडीकाम करणाऱ्याकडून आपली \VUgras{घसरगाडी}
\textsuperscript{(32)} घेतली, तिच्यात आपल्या एका सोबत्याला बसवले आणि धावत निघाला.

%%K t07-c1-17-1 p
१७. --- दुसऱ्यांनी \VUgras{बर्फाच्या ढिगाकडे} \textsuperscript{(37)} धाव घेतली —
\VUgras{पुलाजवळच्या} \textsuperscript{(40)} — आणि कालच उभारलेल्या \VUgras{बर्फाच्या पुतळ्यावर}
\textsuperscript{(39)} मारा सुरू केला; त्याला त्यांनी टोपी, चिलीम आणि खांद्यावरचे दप्तर चढवले
होते.

%%K t07-c1-18-1 p
१८. --- गोठवणारा वारा आणि थंडी यांची त्यांना फिकीर नाही. बहुतेकांच्या गळ्यात \VUgras{मफलरही}
\textsuperscript{(35)} नाही, आणि \VUgras{बर्फाच्या गोळ्यांच्या} \textsuperscript{(38)}
लढाईनंतर पुन्हा ऊब यायला त्यांना ओंजळभर गरमागरम \VUgras{चेस्टनट} \textsuperscript{(42)} पुरतात
— म्हातारी \VUgras{विक्रेतीण} जाकोबिना यांच्याकडून विकत घेतलेले, जी आपल्या फार पातळ
\VUgras{शालीखाली} \textsuperscript{(43)} थंडीने कुडकुडते.

%%K t07-tit-5 sub
\VUsaut{3.0mm}
\VUcentre{12.0pt}{0}{\textit{दुसरा प्रवेश.}}
\VUsaut{3.0mm}

%%K t07-tit-6 sub
\VUpk{10.2pt}{40}{वसंतातले शेतघर.}
\VUsaut{4.0mm}

%%K t07-c2-01-1 p
१. --- हिवाळ्याची सगळी मौज असूनही, \VUgras{वसंताच्या} परतण्याचे आपण मनापासून आनंदाने स्वागत
करतो.

%%K t07-c2-01-2 p
मे महिन्यात, संध्याकाळी, शेताचे अंगण केवढे आनंददायी चित्र दिसते!

%%K t07-c2-02-1 p
२. --- क्षितिजावर \VUgras{सूर्य} \textsuperscript{(1)} हळूहळू मावळतो आहे, \VUgras{शिवार}
\textsuperscript{(3)} आपल्या जांभळ्या \VUgras{किरणांनी} \textsuperscript{(2)} भरून टाकत.
मेहनती \VUgras{नांगरणारा} \textsuperscript{(6)} आपले \VUgras{शेत} \textsuperscript{(4)}
नांगरायची घाई करतो आहे.

%%K t07-c2-02-2 p
आपल्या \VUgras{नांगराने} \textsuperscript{(5)} — धिप्पाड \VUgras{घोड्याला}
\textsuperscript{(7)} जुंपलेल्या — तो \VUgras{तास} \textsuperscript{(8)} पाडतो, ज्यात
\VUgras{पेरणारा} \textsuperscript{(9)} ओंजळभर \VUgras{बियाणे} टाकील, आणि त्यातून सोनेरी कणसे
येतील.

%%K t07-c2-03-1 p
३. --- \VUgras{गवत कापणाऱ्यांनी} \textsuperscript{(11)} आपल्या विळ्यांनी कुरणातले गवत कापले,
वैरण करणाऱ्यांनी \VUgras{वैरण} \textsuperscript{(10)} वाळवली — \VUgras{काट्यांनी}
\textsuperscript{(12)} आणि दंताळ्यांनी उलटीपालटी करत —, आणि हे बघ ते परतत आहेत.

%%K t07-c2-03-2 p
काही जण बैलांनी ओढलेल्या जड गाडीपुढे चालतात; आपले हत्यार ते खांद्यावर घेऊन जातात. दुसरे, जरा
आळशी, सुगंधी वैरणीवर आरामात पसरले आहेत.

%%K t07-c2-04-1 p
४. --- जाताजाता ते \VUgras{माळ्याला} \textsuperscript{(16)} मित्रत्वाने हात करतात, जो
\VUgras{फळबागेत} \textsuperscript{(13)} \VUgras{फळझाडे} छाटण्यात आणि
\VUgras{नासपतीच्या झाडांना} \textsuperscript{(14)} कलम करण्यात गुंतला आहे.
\VUgras{आलुबुखाराच्या झाडांना} \textsuperscript{(15)} बहर येऊ लागला आहे, आणि चांगले खत
घातलेल्या \VUgras{भाजीच्या मळ्यात} \textsuperscript{(17)} भाज्या, कोबी आणि कोशिंबिरीची पाने
आधीच जोरात आहेत.

%%K t07-c2-04-2 p
कुंपणाच्या भिंतीवर एक सुंदर \VUgras{मोर} \textsuperscript{(18)} आपला पिसारा फुलवतो आहे, आपली
देखणी पिसे दाखवण्यासाठी.

%%K t07-tit-7 sub
\VUpk{10.2pt}{40}{शेताचे अंगण.}
\VUsaut{3.0mm}

%%K t07-c2-05-1 p
५. --- \VUgras{तबेल्याची} \textsuperscript{(19)} दारे उघडी आहेत आणि आत गव्हाण आणि
\VUgras{पाचोळ्याचे अंथरूण} \textsuperscript{(23)} दिसते — \VUgras{घोडीचे}
\textsuperscript{(21)}, जिला \VUgras{मोतद्दाराने} \textsuperscript{(20)} नुकतेच जुंपणीतून
सोडले. आपल्या लांब पायांनी इतके गमतीदार दिसणारे \VUgras{शिंगरू} \textsuperscript{(22)} उड्या
मारत आपल्या आईमागे जाते आहे.

%%K t07-c2-06-1 p
६. --- \VUgras{विहिरीजवळ} \textsuperscript{(29)} \VUgras{गायी} \textsuperscript{(25)} लाकडी
\VUgras{डोणीत} \textsuperscript{(26)} पाणी प्यायला जातात, जी त्यांच्यासाठी पाणवठा म्हणून
वापरतात. \VUgras{वासरू} \textsuperscript{(24)} ही संधी साधून दूध पिते, आणि गोठ्यातला चांगला
वास घेत \VUgras{वळू} \textsuperscript{(27)} आनंदाने डरकाळतो आणि दणकट \VUgras{शिंगांचे}
\textsuperscript{(28)} डोके अभिमानाने उंचावतो.

%%K t07-c2-07-1 p
७. --- सगळे कामात आहेत. \VUgras{नोकराणी} \textsuperscript{(32)} विहिरीचा \VUgras{दोर}
\textsuperscript{(31)} हाताने धरते. जनावरांना पाणी पाजण्यासाठी ती \VUgras{बादलीने}
\textsuperscript{(30)} पाणी काढते. \VUgras{गंजीजवळ} \textsuperscript{(33)} एक माणूस पेंढा गोळा
करतो, आणि \VUgras{घराच्या} \textsuperscript{(34)} दारासमोर म्हातारी \VUgras{सूत कातणारी}
\textsuperscript{(35)} आपल्या चरख्याने आणि आपल्या \VUgras{टकळीने} पूर्वीच्या काळासारखे
\VUgras{जवस} किंवा \VUgras{ताग} कातते.

%%K t07-tit-8 sub
\VUsaut{3.0mm}
\VUpk{10.2pt}{40}{शेतकरी.}
\VUsaut{3.0mm}

%%K t07-c2-08-1 p
८. --- \VUgras{शेतकरी} \textsuperscript{(36)} ही सगळी कामे चालवतो आणि आपल्या गड्यांना सूचना
देतो. \VUgras{कुळव} \textsuperscript{(40)} आणि \VUgras{टिकाव} \textsuperscript{(39)} छपराखाली
घ्यायला तो सांगतो, जे \VUgras{घराजवळ} \textsuperscript{(38)} विसरले गेले —
\VUgras{कुत्र्याच्या} \textsuperscript{(37)} घराजवळ.

%%K t07-c2-09-1 p
९. --- त्याच्या हाकांनी एक \VUgras{कबूतर} \textsuperscript{(41)} घाबरते आणि पंख फडफडवत
\VUgras{कबूतरखान्यात} \textsuperscript{(42)} शिरते, आणि \VUgras{कोंबड्या}
\textsuperscript{(43)} \VUgras{खुराड्यात} \textsuperscript{(44)} परतायची घाई करतात.

%%K t07-c2-10-1 p
१०. --- \VUgras{मेंढवाड्यासमोर} \textsuperscript{(48)}, हा बघा म्हातारा \VUgras{धनगर}
\textsuperscript{(45)}, आपला \VUgras{कळप} \textsuperscript{(49)} परत आणणारा. आपली
\VUgras{काठी} \textsuperscript{(46)} हातात धरून — जी त्याच्याजवळ कधीच नसते असे होत नाही, तसाच
त्याचा रंग उडालेला \VUgras{झगाही} \textsuperscript{(47)} — तो गोठ्याकडे \VUgras{मेंढे}
\textsuperscript{(50)}, \VUgras{मेंढ्या} \textsuperscript{(53)}, \VUgras{कोकरे}
\textsuperscript{(54)}, \VUgras{शेळ्या} \textsuperscript{(51)} आणि \VUgras{करडे}
\textsuperscript{(52)} नेतो. त्याचा इमानी \VUgras{कुत्रा} \textsuperscript{(55)} आर्गुस
सगळ्यात शेवटी येतो आणि भुंकून मागे राहणाऱ्यांना घाई करायला लावतो.

%%K t07-c2-11-1 p
११. --- हा सगळा गोंगाट आणि हालचाल त्या भल्या \VUgras{दुभत्या गायीची} \textsuperscript{(56)}
शांतता ढळू देत नाही, जी आपली \VUgras{घंटी} \textsuperscript{(57)} किणकिणवते, तिचे दूध
\VUgras{गोठ्याच्या} \textsuperscript{(58)} दारासमोर काढले जात असताना.

%%K t07-c2-12-1 p
१२. तिला जणू कळते की तिने आपले दूध द्यायला हवे, म्हणजे \VUgras{शेतकरणीला}
\textsuperscript{(60)} \VUgras{रवीत} \textsuperscript{(61)} \VUgras{लोणी} काढता येईल, किंवा ती
उत्तम \VUgras{चीज} \textsuperscript{(64)} करता येतील, जी \VUgras{घंगाळावर}
\textsuperscript{(63)} ठेवलेल्या फळीवरून निथळत असतात.

%%K t07-c2-13-1 p
१३. --- सगळे \VUgras{दूधघर} \textsuperscript{(59)} \VUgras{गडी} \textsuperscript{(65)} फार
स्वच्छ राखतो, ज्याने हातात घेतलेल्या मोठ्या बादलीत नुकतीच \VUgras{दुधाची भांडी}
\textsuperscript{(62)} धुतली.

%%K t07-tit-9 sub
\VUsaut{3.0mm}
\VUpk{10.2pt}{40}{कोंबडवाडा.}
\VUsaut{3.0mm}

%%K t07-c2-14-1 p
१४. --- आपल्या जाड लाकडी वहाणांनी तो जरा अवजड चालतो, आणि त्यामुळे \VUgras{टर्की}
\textsuperscript{(68)}, \VUgras{मादी बदके} \textsuperscript{(66)}, \VUgras{नर बदके}
\textsuperscript{(67)} आणि \VUgras{हंस} \textsuperscript{(69)} पळून जातात, जे
\VUgras{आपली चोच} \textsuperscript{(70)} रुंद उघडून अस्वस्थपणे ओरडतात.

%%K t07-c2-14-2 p
\VUgras{कोंबडा} \textsuperscript{(71)}, जरा अधिक भांडखोर, आपला लाल \VUgras{तुरा} \textsuperscript{(72)} ताठ करतो, आपल्या \VUgras{शेपटीची} \textsuperscript{(73)} पिसे झटकतो आणि « कुकुच्कू » असे दणदणीत बांग देतो.

%%K t07-c2-15-1 p
१५. --- \VUgras{कोंबडी-आई} \textsuperscript{(74)} \VUgras{उकिरड्यावर} \textsuperscript{(81)}
टेहळणी करते, आपल्या \VUgras{पिल्लांसह} \textsuperscript{(75)}. \VUgras{डुकराचे पिल्लू}
\textsuperscript{(78)} आपल्या आईकडे पळते — त्या अवाढव्य \VUgras{डुकरिणीकडे}
\textsuperscript{(77)}, जी आपल्याला डुकरांच्या खुराड्याजवळ दिसते.

%%K t07-c2-16-1 p
१६. --- आपल्या खणांमध्ये \VUgras{ससे} \textsuperscript{(79)} शेतकऱ्याच्या मुलीने आणलेले कोवळे
\VUgras{गवत} \textsuperscript{(80)} हावरेपणाने खातात. जॉनिनाला आपले ससे फार आवडतात, आणि रोज
खुराड्यातून ताजी \VUgras{अंडी} \textsuperscript{(76)} आणून झाली की ती आपल्या लहानग्या
मित्रांना भेटायला येते.

\VUsaut{6.0mm}
\VUfilet{22mm}
